% Student-facing setup flow; included by 01-getting-started.tex.
\subsection{Set up your project with your advisor}
\label{guide:setup:terminal}

Create or reuse your personal GitHub, Overleaf, ORCID, Zotero and Zenodo accounts.
Do not register a second ORCID if you already have one. The setup command displays
the official signup and profile links and asks for non-secret identifiers:
\href{https://github.com/signup}{GitHub},
\href{https://www.overleaf.com/register}{Overleaf},
\href{https://orcid.org/register}{ORCID},
\href{https://www.zotero.org/user/register}{Zotero}, and
\href{https://zenodo.org/signup/}{Zenodo}.
Install Zotero and its browser connector. Use your university-provided identity
for institutional services; conference registration is a separate dated task.

Enable two-factor authentication before accepting project access. The advisor checks the organization requirement before inviting you. An identifier with valid syntax
does not prove account ownership, project access or two-factor authentication.
Enter passwords, recovery codes and API tokens only into the relevant service,
never into the setup identifier prompts.

\begin{enumerate}
  \item The student installs Git, Python 3.11 or newer, GitHub CLI and Docker. On Windows, run
    the setup commands in WSL. On macOS, use the supported Linux-container route.
    Choose the Docker runtime explicitly; Docker Desktop with WSL integration and
    a separate in-distribution engine must not be installed as competing defaults.
  \item The advisor creates the private organization project from the template,
    sets its identity and author/title metadata, and grants the student write
    access to that project. The student accepts the invitation. Repository and
    access automation requires verified account/organization permissions.
  \item The student clones the project and runs
    \texttt{python3 scripts/setup.py student} from its root. Each accepted account
    answer is saved privately before the next prompt. Stop and resume without
    repeating valid answers. Use \texttt{--dry-run} to see the remaining steps
    without changing files or services.
  \item Preview installation with \texttt{--install-plan}. Apply the displayed
    package commands only when intended, using \texttt{--install-tools} and an
    explicit \texttt{--docker-runtime}. Package installation is not proof of a
    working daemon. Verify tools and container execution with
    \texttt{--skip-accounts --check-bootstrap --smoke}. This explicitly downloads
    the selected latest smoke image and executes its resolved digest.
  \item Generate the editor workspace with
    \texttt{python3 scripts/setup.py student --skip-accounts --workspace}.
    Open \texttt{tdk.code-workspace}; existing edited workspace files are preserved.
    Optional local proof and benchmark tools follow the selected components.
  \item The advisor imports the whole GitHub repository into premium Overleaf,
    invites the paper collaborators, configures the root draft/final entrypoints,
    and arranges Zotero group membership and bibliography refresh.
  \item Together, verify compilation and both native synchronization directions.
    Paper collaborators edit the paper in Overleaf and use their normal tools for
    other components. This convention does not protect file modes or symlinks;
    verify code, proof and benchmark execution from a fresh clone after syncing.
\end{enumerate}

Run \texttt{python3 scripts/setup.py status} to see remaining tasks.
The shared \texttt{project.json} contains project identity, components and bounded
metadata. Private account answers and local receipts belong only in the ignored
\texttt{.onboarding/} directory. Do not commit or publish that directory.
The advisor previews changes before \texttt{--apply-metadata}; conflicts preserve
student edits, notes, the chosen lint phase and the Zotero-managed bibliography.

The lead student owns the Zenodo records. Automated packaging prepares a draft;
after advisor review of the exact receipt, the lead student deliberately invokes
the project publication workflow with their own Zenodo repository secret.
Review is a working policy, not a technical approval boundary on student-writable
workflows. Account registration and a completed local checklist do not authorize
publication or prove remote integration.
